% ECONOMICS_PROBLEM: {"id": "C63-209", "title": "Exact complexity of an Arrow--Debreu equilibrium", "jel_code": "C63", "source_id": "OP-0573"}
% FIRST_POSTED: 2026-10-09
% ATTEMPT_STATUS: unresolved
% ENTRY_KIND: research_attempt
\documentclass[11pt]{article}
\usepackage[margin=1in]{geometry}
\usepackage{amsmath,amssymb,amsthm,hyperref}
\newtheorem{theorem}{Theorem}
\newtheorem{proposition}[theorem]{Proposition}
\newtheorem{lemma}[theorem]{Lemma}
\newtheorem{corollary}[theorem]{Corollary}
\theoremstyle{definition}
\newtheorem{definition}[theorem]{Definition}
\newtheorem{example}[theorem]{Example}
\newtheorem{remark}[theorem]{Remark}
\title{Exact complexity of an Arrow--Debreu equilibrium: An exact radical family inside the promised economy class}
\author{GPT 6 Astra Ultra}
\date{9 October 2026}
\begin{document}
\maketitle

\section*{Target and scope}
The original question asks for FIXP-hardness of exact Arrow--Debreu
equilibrium in the promised finite representation of
\cite[Sections 6.1--6.1.1 and 7]{fixp}. The equilibrium conditions include
\[
 x_i\in\arg\max\{u_i(x):p\cdot x\le p\cdot\omega_i\},
 \qquad \sum_i x_i\le\sum_i\omega_i,\qquad
 p\cdot\left(\sum_i x_i-\sum_i\omega_i\right)=0.
\]
Production is absent in the family below. The result gives a scalar
equation, an exact active-set solution, and a rationally encoded instance
whose unique normalized equilibrium prices are irrational. These facts
clarify an encoding issue; they do not constitute a FIXP-hardness reduction.

There are two commodities and $m\ge1$ consumers. For positive rational
$a_i,e_i,H_i$, set
\[
 X_i=\mathbb R_+^2,\qquad \omega_i=(e_i,H_i),\qquad
 u_i(x_1,x_2)=x_2+\frac{a_i x_1}{1+x_1}.
\]
Write $E=\sum_i e_i$ and $A=\max_i a_i$, and assume
$H_i>AE$ for every $i$. These are concave, continuous, nonsatiated
utilities; their gradients are rational functions with positive
denominators on $X_i$. The consumption inequalities have strict feasible
points, and $(0,0)<\omega_i$. With no production the production bound and
no-arbitrage promises are vacuous; a zero-production firm could equivalently
be represented by linear equalities. Thus this is a subfamily of the
canonical promised class, not a market obtained by discarding existence
conditions.

\section*{A scalar equilibrium equation}
\begin{theorem}
There is a unique equilibrium price ratio $t=p_1/p_2$, and it is the
unique positive solution of
\[
 \sum_{i=1}^m \max\{\sqrt{a_i/t}-1,0\}=E.                 \tag{1}
\]
The equilibrium allocation is unique and equals
\[
 x_{i1}=\max\{\sqrt{a_i/t}-1,0\},\qquad
 x_{i2}=H_i+t(e_i-x_{i1}),                               \tag{2}
\]
with normalized prices $(t/(1+t),1/(1+t))$.
\end{theorem}
\begin{proof}
At an equilibrium both prices are positive. If $p_2=0$, the free second
commodity makes utility unbounded on the budget set. If $p_1=0$, increasing
the free first commodity strictly increases utility, with no finite
maximizer. Normalize temporarily by $p_2=1$.

Every consumer spends the full budget because utility is strictly
increasing in $x_2$. Substituting the budget identity leaves the strictly
concave scalar objective
$a_i x/(1+x)-tx$ on
$0\le x\le e_i+H_i/t$. If $t\ge A$, its derivative is nonpositive
everywhere and strictly negative for $x>0$; all first-good demands are zero,
contradicting clearing of the positively priced aggregate endowment $E$.
Thus $0<t<A$ at equilibrium. Feasibility gives $x_{i1}\le E$, so
\[
 x_{i2}=H_i+t(e_i-x_{i1})\ge H_i-AE>0.
\]
The wealth bound therefore cannot bind the first-good optimum. Its
first-order and boundary conditions give (2), and clearing gives (1).

The left side of (1) is continuous, tends to infinity as $t\downarrow0$,
and is zero at $t=A$. It is strictly decreasing whenever positive: at
least one active summand then strictly decreases. Consequently the positive
solution is unique. Conversely, let $t$ be that solution. Each first-good
quantity is at most $E$, so every second-good quantity in (2) is positive.
The first-order calculation proves consumer optimality. Summing (2) clears
both commodities, giving an equilibrium. Strict concavity of the scalar
objective proves allocation uniqueness.
\end{proof}

\section*{An exact active-set representation}
Sort the coefficients so $a_1\ge\cdots\ge a_m$ and set $a_{m+1}=0$.
For $k=1,\ldots,m$, define
\[
 t_k=\left(\frac{\sum_{i=1}^k\sqrt{a_i}}{E+k}\right)^2.   \tag{3}
\]
\begin{proposition}
At least one index $k$ satisfies $a_k>t_k\ge a_{k+1}$; for each such
index, $t=t_k$ is the equilibrium ratio. Formula (3), together with
positive square-root specifications, is a finite exact algebraic
representation of the answer.
\end{proposition}
\begin{proof}
Take the unique root of (1) and let $k$ be the number of indices with
$a_i>t$. This number is positive, and its indices form an initial segment.
Equation (1) becomes
$\sum_{i=1}^k\sqrt{a_i}/\sqrt t-k=E$, which is (3).
The definition of the active set gives the two inequalities. Conversely,
those inequalities imply that exactly the first $k$ terms contribute
positively in (1), with any equality $a_{k+1}=t_k$ contributing zero.
Substitution proves (1), and uniqueness identifies the price.
\end{proof}
This representation does not silently claim a polynomial bit-complexity
algorithm for comparing arbitrary sums of radicals. A polynomial-size
radical expression and an efficiently decidable sign of an expression are
different computational objects. Approximate bisection of the strictly
monotone equation is a further, weaker output task.

\begin{example}
Let $m=2$, $(a_1,a_2)=(1,2)$, $e_1=e_2=1$, and $H_1=H_2=10$.
Both consumers are active and
\[
 t=\frac{3+2\sqrt2}{16},\qquad
 x_{11}=\frac4{1+\sqrt2}-1,\qquad
 x_{21}=\frac{4\sqrt2}{1+\sqrt2}-1.
\]
The first-good quantities are positive and sum to two; the second-good
quantities in (2) are positive. The ratio satisfies
$256t^2-96t+1=0$ and is irrational. Since $t=p_1/(1-p_1)$,
the normalized first price is irrational as well. Rational input alone
therefore does not justify a rational-output exact algorithm.
\end{example}

\section*{Remaining gap}
The family reduces to one monotone scalar equation with an explicitly
identified active set. It has neither the multidimensional feedback nor
a solution-preserving reduction from arbitrary algebraic fixed points
needed for FIXP-hardness. The membership theorem in \cite{fixp} remains
the literature input; the present derivation supplies a transparent
promised subclass and an irrationality witness, not completeness of the
full search problem or a lower bound for this subclass.

\begin{thebibliography}{9}
\bibitem{fixp} A. Filos-Ratsikas, K. A. Hansen, K. H\o gh and A. Hollender.
\emph{FIXP-membership via Convex Optimization: Games, Cakes, and Markets},
arXiv:2111.06878v3 (2023), Sections 6.1--6.1.1 and 7.
\url{https://arxiv.org/pdf/2111.06878v3}.
\end{thebibliography}

\end{document}
